
% ════════════════════════════════════════════
% 模型的改进与推广（08.模型改进推广.tex）—— 写作提示词见本文件末尾【本章节写作提示词】。
% ════════════════════════════════════════════
\section{模型的改进与推广}

\subsection{模型的改进}

本文模型可通过以下方向进一步改进。在方法层面，可将固定网格枚举替换为基于梯度或序列二次规划的连续优化，消除离散误差并支持更大规模站点选择；在多目标层面，可将定位精度指标与最大观测开始时间加权为统一目标，得到速度与精度的帕累托前沿；在参数层面，可将行驶速度、云台角速度视为随机变量，引入鲁棒优化或机会约束，使调度方案对参数扰动不敏感；在数据层面，可结合真实地形与动态障碍信息，将建筑物从轴对齐长方体推广为任意多边形棱柱，并加入车辆避碰与转向时间，使时间轴更贴近实际。

\subsection{模型的推广}

本文建立的模型框架具有较强的通用性，可推广到更广泛的协同探测与调度场景：其一，可推广到多无人机目标跟踪与接力定位，将静止目标假设扩展为运动目标，并引入预测与滤波；其二，可推广到城市安防中多传感器布站与应急响应调度，如地面摄像头、声学传感器与巡逻车的协同部署；其三，可推广到物流运输与应急车辆路径规划，在突发任务与设备故障下实现动态重调度；其四，误差传播与精度约束方法可推广到任何基于角度交会的定位系统，如无线电测向、声呐定位与天文导航等，也可与无线传感器网络节点定位算法相融合\cite{yang2009dingwei}。
